\documentclass[instructions.tex]{subfiles}
\begin{document}
 
\chapter{Illusion au sujet de la bonne mort.}
 
Chacun désire de bien mourir ; mais les uns ne savent ce que c’est qu’une bonne mort, et les autres se mettent peu en peine de s’y préparer. Illusion dans les uns et dans les autres.
 
\primo Tous ceux qui meurent avec des signes de repentir, avec les sacremens, les larmes aux yeux, le crucifix en main, n’ont pas toujours le bonheur de mourir de la mort des justes. Ces marques sont équivoques, et le cœur n’a peut-être point de part à toute cette piété apparente.
 
Faire une bonne mort, c’est mourir sans péché mortel, dans l’amour de Dieu et dans sa grâce ; voilà ce que vous devez désirer. Quand vous mourriez de la mort la plus déplorable aux yeux des hommes, pourvu que vous mouriez dans la grâce de Dieu, votre salut est assuré ; au contraire, quand vous auriez donné toutes les marques de pénitence, et reçu tous les sacremens, si vous mouriez dans l’affection volontaire à un seul péché mortel, et sans amour de Dieu, il n’y aurait point de salut pour vous. Les sacremens sont plus utiles et plus nécessaires à la mort ; mais ils sont infructueux, et ne procurent pas une sainte mort, si le cœur n’est pas converti. Erreur de penser autrement.
 
\secundo C’est encore une illusion qui n’est pas moins funeste, de s’attendre à une sainte mort, sans faire pendant la vie ce qu’il faut pour bien mourir. Vous voudriez mourir saintement ; vivez donc comme vous voudriez mourir. La mort est l’écho de la vie ; une vie sainte est suivie d’une sainte mort, comme la mauvaise mort est la suite ordinaire d’une mauvaise vie.
 
Ne vous trompez pas, dit saint Paul, on ne se moque pas de Dieu en vain. L’homme ne recueillera que ce qu’il aura semé. Après avoir semé de l’ivraie, vous ne devez pas attendre de bon grain ; vous devez donc attendre une fin malheureuse, si vous vivez mal \footnote{Quæ enim seminaverit homo, hæc et metet. Gal. 6. 8.}.
 
C’est une témérité, en vivant dans le crime, de s’attendre à une mort sainte. On ne trouve dans l’Ecriture, dit saint Bernard, qu’un homme qui soit mort saintement après avoir mal vécu (c’est le bon larron). Il s’en trouve un, pour ôter tout sujet de désespoir au pécheur ; mais il ne s’en trouve qu’un seul, pour lui ôter tout sujet de présomption .\footnote{Unus ne desperes; unicus, ne præsumas. S. Bern.}.
 
La conduite ordinaire de la Providence, c’est de laisser aller les choses selon l’ordre qu’elle a établi. Dieu laisse agir le feu selon son activité ; il laisse couler les fleuves selon leur pente : de même il laisse mourir en impénitens ceux qui ont vécu dans l’impénitence. Quand il arrive autrement, c’est un miracle ; Dieu fait quelquefois ce miracle pour des raisons que nous devons adorer ; mais c’est une présomption de s’y attendre en vivant mal.
 
Oh ! si vous compreniez les horreurs d’une mauvaise mort, quelles précautions ne prendriez-vous pas pour vous en préserver ! Mourir dans le péché ! mourir ennemi de Dieu ! mourir en réprouvé ! Quelle mort ! Oh ! qu’on est aveugle, si on ne travaille pas à éviter un si grand malheur !
 
Si vous compreniez, au contraire, les douceurs d’une mort, vous comprendriez qu’un moment de ces consolations vous dédommagerait de toutes les peines que vous auriez au service de Dieu. Qu’il est doux de dire alors comme David : Dans les ombres de la mort je n’ai rien à craindre, parce que j’ai mon Dieu avec moi \footnote{In medio umbræ mortis non timebo mala, quoniam tu mecum es. Ps. 22. 4.}. Tels furent les sentimens d’un pieux personnage, qui, ayant reçu le sacré viatique, sentit dans son âme tant de délices, qu’il s’écria : Je ne m’attendais pas à mourir avec tant de douceur et de joie\footnote{ Non putabam tam jucundè mori.}.
 
\section{Résolutions}
 
\primo Conversion sincère et entière dès aujourd’hui. 

\secundo Prière fréquente pour demander à Dieu la grâce d’une sainte mort. --

\tertio Et vie désormais telle que je voudrais l’avoir passée à l’heure de mon trépas. 

\prayer{Faites, ô mon Dieu ! que je vive dès maintenant comme je désirerai avoir vécu, quand il me faudra quitter la vie présente pour aller à votre jugement.}
 
\end{document}
